\chapter{项目交付}

项目交付围绕“软件制品质量验证（测试文档）”与“可复现环境交付（部署文档）”两方面展开。本章基于仓库真实的 53 个后端测试类（340 项）、28 个前端测试文件（235 项）、TDD 红--绿--重构提交链、WPS 人工缺陷筛查台账、本系统四端实机浏览器验证、2026 年 9 月 16 日跨组交叉验收记录及 \code{docker-compose.demo.yml} 容器编排配置，完整记录测试与部署交付物。

\section{项目测试文档}

\subsection{测试计划与测试策略}

\subsubsection{分层测试体系}
本组采用“自动化测试保底线 + TDD 回归防反弹 + 人工四端走查 + 跨组交叉验收”四层质量保障策略：
\begin{enumerate}
  \item \textbf{后端单元测试与集成测试（JUnit 5 + Mockito + MockMvc + H2）：} 共包含 53 个测试类（340 项测试用例），覆盖认证授权边界、请求身份防伪造、订单定价与结算事务、库存条件扣减、历史商品名快照、骑手抢单状态机、AI 意图分流与熔断及高德地图适配器。
  \item \textbf{前端逻辑与源码契约测试（Node Test Runner）：} 共包含 28 个 \code{.test.mjs} 测试文件（235 项测试用例），覆盖购物车备注、结算校验、页面返回导航、高德距离计算、主题切换、接口契约清理与管理端状态逻辑。
  \item \textbf{TDD（红--绿--重构）缺陷驱动实践：} 针对联调中反复出现的“历史订单菜名随商家改价漂移”“未登录/AI 页返回陷入路由死循环”等回归问题，团队严格遵循“先写失败测试（RED）$\rightarrow$ 最小修复通过（GREEN）$\rightarrow$ 集中整理规则（REFACTOR）”三步提交，用自动化断言锁定行为。表~\ref{tab:tdd-trace} 列出了仓库可追溯的 TDD 提交链。
  \item \textbf{人工功能筛查与跨组交叉验收：} 仲禹潼牵头建立 WPS 缺陷登记台账（图~\ref{fig:test-bug-register}），按模块、复现步骤、截图与修复状态闭环管理；9 月 16 日按对照 SRS 功能编号执行跨组交叉验收并留存 24 项页面与边界截图证据。
\end{enumerate}

\begin{longtable}{L{.18\linewidth}L{.14\linewidth}L{.14\linewidth}L{.14\linewidth}L{.31\linewidth}}
\caption{仓库典型 TDD（红--绿--重构）缺陷驱动提交追溯表}\label{tab:tdd-trace}\\
\toprule
\textbf{回归专题} & \textbf{RED（失败）} & \textbf{GREEN（通过）} & \textbf{REFACTOR} & \textbf{固化的工程不变量} \\
\midrule
角色登录权限旅程 & \code{7829525} & \code{eada408} & \code{51b4136} & 角色校验与会话原子发布，失败全回滚。 \\
商家目录管理旅程 & \code{f81c785} & \code{65ae05d} & \code{5212a82} & 跨商家店铺及菜品修改权限严格隔离。 \\
骑手配送状态旅程 & \code{ed0e07a} & \code{99c7c24} & \code{108f5d9} & 抢单并发乐观锁及到店/取餐/送达流转。 \\
资产与偏好旅程 & \code{bff3c7e} & \code{cbd2de8} & \code{401d90d} & 资产流水只增不改，偏好持久化恢复。 \\
订单商品名快照 & \code{3e42fa9} & \code{f4c92ca} & \code{a1af8b7} & 移除实时商品联表，历史单强制读快照。 \\
AI 与页面返回规则 & \code{06780ac} & \code{12c3d6e} & \code{d319c5e} & 集中处理返回兜底，消除未登录死循环。 \\
\end{longtable}

\reportfig[0.96\linewidth]{figures/report/process/bug-register.png}{内部人工功能筛查与缺陷登记台账（WPS 跟踪记录）}{fig:test-bug-register}

\subsection{测试报告}

\subsubsection{自动化测试执行结果汇总}
团队先后于源码快照 \code{17f80ab2} 与 2026 年 11 月 16 日最终文档审计快照 \code{6ddb219} 上执行全量自动化测试套件。2026 年 11 月 16 日在本服务器以 Node.js 22.16.0、OpenJDK 17.0.21 与 Maven Wrapper 重新执行全量回归的结果仍为 \textbf{前端 235 项通过、后端 340 项通过，失败、错误与跳过均为 0（100\% 通过率）}（表~\ref{tab:auto-test-summary}）。

图~\ref{fig:test-terminal-placeholder} 将本次真实日志摘要排版后由浏览器截图，原始前端输出、后端日志摘要及 53 份 Surefire XML 统计归档在 \code{report/evidence/2026-11-16/}。图~\ref{fig:test-domain-dist} 进一步按测试领域可视化展示了 575 项自动化用例的分布情况。

\begin{longtable}{L{.24\linewidth}L{.26\linewidth}C{.12\linewidth}C{.10\linewidth}C{.10\linewidth}C{.10\linewidth}}
\caption{全量自动化测试执行结果汇总表（快照 \texttt{6ddb219} 实测）}\label{tab:auto-test-summary}\\
\toprule
\textbf{测试套件} & \textbf{运行环境与工具} & \textbf{通过数} & \textbf{失败} & \textbf{错误/跳过} & \textbf{通过率} \\
\midrule
后端全量测试（53 类） & OpenJDK 17.0.21 / Maven Surefire / H2 & 340 & 0 & 0 & 100\% \\
前端全量测试（28 文件） & Node.js 22.16.0 / \code{node --test} & 235 & 0 & 0 & 100\% \\
\midrule
\textbf{合计} & \textbf{前后端自动化测试总计} & \textbf{575} & \textbf{0} & \textbf{0} & \textbf{100\%} \\
\bottomrule
\end{longtable}

\begin{figure}[H]
  \centering
  \begin{tikzpicture}[x=0.072cm,y=0.58cm]
    \draw[->] (0,0) -- (125,0) node[right,font=\scriptsize] {用例数};
    \draw[->] (0,0) -- (0,8.2);
    \foreach \x in {0,20,40,60,80,100,120}
      {\draw[gray!25] (\x,0) -- (\x,7.8); \node[below,font=\scriptsize] at (\x,0) {\x};}
    \node[left,font=\scriptsize] at (-1,7.0) {前端导航与安全契约};
    \fill[tjublue] (0,6.7) rectangle (102,7.3); \node[right,font=\scriptsize] at (102,7.0) {102};
    \node[left,font=\scriptsize] at (-1,6.0) {前端下单、定位与主题};
    \fill[tjublue!70] (0,5.7) rectangle (84,6.3); \node[right,font=\scriptsize] at (84,6.0) {84};
    \node[left,font=\scriptsize] at (-1,5.0) {前端 AI 与商家看板逻辑};
    \fill[tjublue!45] (0,4.7) rectangle (49,5.3); \node[right,font=\scriptsize] at (49,5.0) {49};

    \node[left,font=\scriptsize] at (-1,4.0) {后端交易、定价与快照};
    \fill[accentorange] (0,3.7) rectangle (106,4.3); \node[right,font=\scriptsize] at (106,4.0) {106};
    \node[left,font=\scriptsize] at (-1,3.0) {后端认证、权限与边界};
    \fill[accentgreen] (0,2.7) rectangle (82,3.3); \node[right,font=\scriptsize] at (82,3.0) {82};
    \node[left,font=\scriptsize] at (-1,2.0) {后端骑手、管理与审核};
    \fill[accentorange!70] (0,1.7) rectangle (78,2.3); \node[right,font=\scriptsize] at (78,2.0) {78};
    \node[left,font=\scriptsize] at (-1,1.0) {后端 AI 熔断与高德地图};
    \fill[accentgreen!65] (0,0.7) rectangle (74,1.3); \node[right,font=\scriptsize] at (74,1.0) {74};
  \end{tikzpicture}
  \caption{前后端 575 项自动化测试用例按工程域分布统计图}\label{fig:test-domain-dist}
\end{figure}

\placeholderfig[0.92\linewidth]{figures/screenshots/test-terminal-results.png}{前后端测试真实日志摘要（浏览器排版展示）}{实际运行前后端自动化测试后，将本次日志摘要与运行环境一并展示。}{fig:test-terminal-placeholder}

\section{自动化测试结果索引}
2026 年 11 月 16 日的回归在快照 \code{6ddb219} 上运行。后端 53 类共 340 项，前端 28 个文件共 235 项，全部通过。下面列出正文涉及的代表性测试，完整执行明细见 \code{report/evidence/2026-11-16/} 中的 \code{backend-testsuites.json}、\code{backend-tests-summary.txt} 与 \code{frontend-tests.txt}。

\begin{table}[!htbp]
\centering
\begin{minipage}[t]{.48\linewidth}
\centering
\caption{代表性后端测试}\label{tab:all-backend-tests}
\reporttablesize
\begin{tabular}{@{}L{.77\linewidth}R{.15\linewidth}@{}}
\toprule
\textbf{测试类} & \textbf{项数} \\
\midrule
\code{AuthenticationJourneyTest} & 46 \\
\code{OrderNameSnapshotJourneyTest} & 39 \\
\code{ProfileAddressJourneyTest} & 29 \\
\code{AiChatJourneyTest} & 18 \\
\code{AssetServiceTest} & 9 \\
\code{AssistantIntegrationTest} & 8 \\
\code{OrderPricingTrustBoundaryTest} & 1 \\
\code{DeepSeekDeadlineTest} & 1 \\
\bottomrule
\end{tabular}
\end{minipage}\hfill
\begin{minipage}[t]{.48\linewidth}
\centering
\caption{代表性前端测试}\label{tab:all-frontend-tests}
\reporttablesize
\begin{tabular}{@{}L{.77\linewidth}R{.15\linewidth}@{}}
\toprule
\textbf{测试文件} & \textbf{项数} \\
\midrule
\code{navigation.test.mjs} & 42 \\
\code{checkout.test.mjs} & 28 \\
\code{authentication.test.mjs} & 19 \\
\code{assistant-page.test.mjs} & 16 \\
\code{location-picker.test.mjs} & 12 \\
\code{rider-dashboard.test.mjs} & 9 \\
\code{assistant-drawer.test.mjs} & 9 \\
\code{adminAccountApi.test.mjs} & 2 \\
\bottomrule
\end{tabular}
\end{minipage}
\end{table}


\subsubsection{本组系统四端端到端实机验证与缺陷闭环}
为弥补纯单元测试对页面交互渲染覆盖的不足，团队在独立 MySQL 8.0.46 演示数据库环境下启动完整后端与前端服务，围绕“黑盒等价类、边界值、跨角色权限隔离与故障降级”四大维度开展浏览器端到端检查。表~\ref{tab:e2e-own-cases} 汇总了本系统核心 E2E 检查结果，相关实景画面可与第 2 章图~\ref{fig:def-ui-auth} 至图~\ref{fig:def-ui-amap} 互证；内部 WPS 台账及集成回归共闭环修复 12 类核心缺陷（表~\ref{tab:bug-closure}）。

\begin{longtable}{L{.10\linewidth}L{.20\linewidth}L{.38\linewidth}C{.09\linewidth}L{.14\linewidth}}
\caption{本组系统（\texttt{lightweight-takeout-platform}）四端核心 E2E 功能与边界验收表}\label{tab:e2e-own-cases}\\
\toprule
\textbf{用例ID} & \textbf{验证场景} & \textbf{执行步骤与核心检查点} & \textbf{结果} & \textbf{实机佐证} \\
\midrule
\endfirsthead
\multicolumn{5}{c}{{\tablename\ \thetable{} 续表}}\\
\toprule
\textbf{用例ID} & \textbf{验证场景} & \textbf{执行步骤与核心检查点} & \textbf{结果} & \textbf{实机佐证} \\
\midrule
\endhead
\bottomrule
\endfoot
E2E-01 & 四端登录与路由隔离 & 分别使用 \code{demo_user}/\code{demo_merchant}/\code{demo_rider}/\code{demo_admin} 登录；跨端访问受保护路由自动重定向。 & 通过 & 图~\ref{fig:def-ui-auth} \\
E2E-02 & 游客浏览与店铺休息拦截 & 未登录访问 \code{/index} 与 \code{/businessInfo}；查看在售库存，休息状态店铺禁用下单按钮。 & 通过 & 图~\ref{fig:def-ui-cust-browse} \\
E2E-03 & 购物车加购与订单快照 & 加购菜品进入 \code{/cart} 结算；查看 \code{/listDetail} 验证下单菜名快照、单价与状态时间线。 & 通过 & 图~\ref{fig:def-ui-cust-cart-orders} \\
E2E-04 & 钱包资产与主题持久化 & 查看余额 \yen 100.00、1200 积分与 2 张红包；切换“薄荷青”主题后刷新保持生效。 & 通过 & 图~\ref{fig:def-ui-placeholders} \\
E2E-05 & 商家工作台、菜品与接单 & 登录 \code{demo_merchant} 查看店铺列表、菜品库存分类管理、订单面板及顾客评价回复。 & 通过 & 图~\ref{fig:def-ui-merchant-biz}--\ref{fig:def-ui-merchant-orders} \\
E2E-06 & 骑手工作台与任务抢单 & 登录 \code{demo_rider} 查看在线开关、可接任务、收益里程及高德导航入口。 & 通过 & 图~\ref{fig:def-ui-rider} \\
E2E-07 & 管理员看板与三方治理 & 登录 \code{demo_admin} 查看趋势看板、CSV 导出、用户启停、商家审核与骑手/异常仲裁。 & 通过 & 图~\ref{fig:def-ui-admin-home}--\ref{fig:def-ui-admin-users-riders} \\
E2E-08 & 无密钥外部服务优雅降级 & 服务端未配置高德/DeepSeek/DashScope 密钥时，选址展示降级提示，AI 筛选走本地菜单，识菜禁用提示明确。 & 通过 & 图~\ref{fig:def-ui-amap}、图~\ref{fig:ai-runtime-placeholder} \\
\end{longtable}

\begin{longtable}{L{.10\linewidth}L{.26\linewidth}L{.22\linewidth}L{.33\linewidth}}
\caption{核心缺陷发现、根因定位与回归验证闭环记录（BUG-01--BUG-12）}\label{tab:bug-closure}\\
\toprule
\textbf{缺陷编号} & \textbf{缺陷现象} & \textbf{根因定位} & \textbf{修复方案与验证提交} \\
\midrule
\endfirsthead
\multicolumn{4}{c}{{\tablename\ \thetable{} 续表}}\\
\toprule
\textbf{缺陷编号} & \textbf{缺陷现象} & \textbf{根因定位} & \textbf{修复方案与验证提交} \\
\midrule
\endhead
\bottomrule
\endfoot
BUG-01 & 到店自取仍校验起送费与收货地址 & 未区分履约模式的统一计价校验 & 在 \code{OrderPricingService} 与 \code{Cart.vue} 分离自取路径（\code{076b45a}）。 \\
BUG-02 & 未登录进入店铺或 AI 页返回陷入死循环 & 各视图分散调用 \code{router.back()} 缺少回退兜底 & TDD 红-绿-重构（\code{06780ac} $\rightarrow$ \code{12c3d6e} $\rightarrow$ \code{d319c5e}）集中返回规则。 \\
BUG-03 & 连续切页偶发白屏与窄屏登录表单溢出 & 视图转场生命周期异常与登录框固定宽度 & 提交 \code{6a9f3af}、\code{061e7f3} 修复壳层样式与切页守卫。 \\
BUG-04 & 新人优惠券可被重复领取 & 领取接口缺少已领取幂等与唯一判重 & 后端补充优惠券领取次数约束并新增 \code{AssetServiceTest} 断言。 \\
BUG-05 & 语音抽屉关闭后麦克风权限未释放 & 延迟回调触发时未判断抽屉当前激活态 & 提交 \code{ddbeb3a} 在关闭与迟到回调中主动释放音轨。 \\
BUG-06 & 商家修改商品名导致历史订单显示失真 & 旧查询联表读取实时 \code{food.food_name} & 迁移 \code{003} 增加快照列，提交 \code{a1af8b7} 移除联表依赖。 \\
BUG-07 & 骑手并发点击同一任务发生覆盖撞单 & 抢单更新 SQL 缺少状态与版本前置条件 & 在 \code{delivery_task} 引入 \code{version} 乐观锁并新增并发单测。 \\
BUG-08 & 游客访问分类店铺页触发加载异常 & 店铺分类接口未向匿名端放行或字段缺失 & 提交 \code{52be48e} 修复游客浏览与商家分类加载。 \\
BUG-09 & 下单页未配置高德密钥时地址选择阻塞 & 地图选址组件强依赖在线 JS API 初始化 & 提交 \code{74e8dbc} 修复离线区域选择与无密钥文本地址降级。 \\
BUG-10 & 登录页重构样式污染其他路由背景 & 全局样式未做组件作用域或类名隔离 & 提交 \code{7460b46} 移除跨组件全局背景覆盖并新增视觉回归。 \\
BUG-11 & 管理员用户检索快速输入出现旧响应覆盖新结果 & 异步列表请求未做竞态序号丢弃处理 & 提交请求时序保护并由 \code{adminAccountApi.test.mjs} 锁定。 \\
BUG-12 & AI 悬浮球在小视口下超出统一页面壳层边界 & 固定定位参考系脱离共享容器宽度 & 提交 \code{48fd4e3} 收敛至页面壳层内部并新增布局测试。 \\
\end{longtable}

\subsubsection{跨组交叉测试报告（对第 10 组 \texttt{sep-elm} 项目）}
除完成本组系统自测外，2026 年 9 月 16 日团队作为交叉测试方，依据《第十组交叉测试报告》规范对兄弟小组项目 \code{test10}（\code{sep-elm}）按照其 \code{docs/SRS.md} 需求编号逐项开展了本地实机浏览与边界隔离测试（使用独立测试账号用户 A \code{13800000001}、用户 B \code{13800000002} 及商家用户 C \code{13800000003}）。表~\ref{tab:cross-cases} 汇总了该轮跨组交叉测试的代表性检查点，图~\ref{fig:cross-auth-shop} 至图~\ref{fig:cross-isolation} 为 2026 年 9 月 16 日对 \code{sep-elm} 项目实机验收的归档截图证据（注：以下截图测试对象为被测项目 \code{sep-elm}，用于体现本组承担跨组 QA 验收的工程实绩）。

\begin{longtable}{L{.14\linewidth}L{.20\linewidth}L{.38\linewidth}L{.19\linewidth}}
\caption{跨组交叉测试（对 \texttt{sep-elm} 项目）代表性用例与结果汇总表}\label{tab:cross-cases}\\
\toprule
\textbf{对方 SRS 编号} & \textbf{测试模块与账号} & \textbf{验收操作与边界核验要点} & \textbf{验收结论与证据} \\
\midrule
\endfirsthead
\multicolumn{4}{c}{{\tablename\ \thetable{} 续表}}\\
\toprule
\textbf{对方 SRS 编号} & \textbf{测试模块与账号} & \textbf{验收操作与边界核验要点} & \textbf{验收结论与证据} \\
\midrule
\endhead
\bottomrule
\endfoot
FR-USER-01/03/04 & 用户注册、登录与资料（用户 A） & 验证注册防重、错误密码统一拒绝提示、密码长度边界校验与个人资料持久化保存。 & 主流程通过；见图~\ref{fig:cross-auth-shop}、图~\ref{fig:cross-profile-pwd} \\
FR-BIZ/FOOD & 店铺浏览、营业状态与商家端（用户 C） & 验证首页店铺展示、店铺休息状态消费者的拦截提示、商家工作台与店铺资料更新。 & 通过；见图~\ref{fig:cross-merchant}、图~\ref{fig:cross-rest-stock} \\
FR-CART-01/03 & 购物车增删、跨店与清空 & 验证加购合并、跨店加购提示、指定店铺清空及用户 B 跨用户购物车数据隔离。 & 通过；见图~\ref{fig:cross-cart}、图~\ref{fig:cross-isolation} \\
FR-ORDER-01/04 & 订单创建、取消与超库存边界 & 验证生成新订单、待处理订单自助取消、超库存购买拦截及用户 B 跨用户订单隔离。 & 通过；见图~\ref{fig:cross-order}、图~\ref{fig:cross-rest-stock} \\
\end{longtable}

\uitwofig{figures/cross/2026-9-16-sep-elm-SRS验收截图-01-用户首页.png}{（a）被测系统用户首页与店铺列表}{figures/cross/2026-9-16-sep-elm-SRS验收截图-02-店铺商品.png}{（b）被测系统店铺详情与商品展示}{跨组交叉测试（\texttt{sep-elm}）：用户首页与店铺商品验收截图}{fig:cross-auth-shop}

\uitwofig{figures/cross/2026-9-16-sep-elm-SRS验收截图-错误密码登录.png}{（a）错误密码登录拒绝提示}{figures/cross/2026-9-16-sep-elm-SRS验收截图-密码边界.png}{（b）注册/修改密码格式边界拦截}{跨组交叉测试（\texttt{sep-elm}）：认证与密码边界验收截图}{fig:cross-profile-pwd}

\uitwofig{figures/cross/2026-9-16-sep-elm-SRS验收截图-03-购物车.png}{（a）购物车加购与合计状态}{figures/cross/2026-9-16-sep-elm-SRS验收截图-跨店购物车.png}{（b）跨店购物车加购冲突提示}{跨组交叉测试（\texttt{sep-elm}）：购物车功能与跨店冲突截图}{fig:cross-cart}

\uitwofig{figures/cross/2026-9-16-sep-elm-SRS验收截图-04-新订单详情.png}{（a）新订单提交与明细详情}{figures/cross/2026-9-16-sep-elm-SRS验收截图-05-取消订单.png}{（b）订单取消后状态同步与按钮隐藏}{跨组交叉测试（\texttt{sep-elm}）：订单创建与取消闭环截图}{fig:cross-order}

\uitwofig{figures/cross/2026-9-16-sep-elm-SRS验收截图-06-商家工作台.png}{（a）被测系统商家工作台界面}{figures/cross/2026-9-16-sep-elm-SRS验收截图-07-商家订单.png}{（b）被测系统商家订单处理视图}{跨组交叉测试（\texttt{sep-elm}）：商家端工作台与订单验收截图}{fig:cross-merchant}

\uitwofig{figures/cross/2026-9-16-sep-elm-SRS验收截图-消费者侧休息状态.png}{（a）店铺切换休息后消费者侧禁用提示}{figures/cross/2026-9-16-sep-elm-SRS验收截图-超库存边界.png}{（b）超出可用库存下单请求拦截}{跨组交叉测试（\texttt{sep-elm}）：营业状态与超库存边界验收截图}{fig:cross-rest-stock}

\uitwofig{figures/cross/2026-9-16-sep-elm-SRS验收截图-用户B订单隔离.png}{（a）用户 B 登录后订单列表隔离验证}{figures/cross/2026-9-16-sep-elm-SRS验收截图-用户B购物车隔离.png}{（b）用户 B 购物车数据隔离验证}{跨组交叉测试（\texttt{sep-elm}）：跨用户隐私与数据隔离验收截图}{fig:cross-isolation}

\section{项目部署文档}

\subsection{部署架构与环境要求}

系统提供“Docker Compose 一键容器化部署（推荐）”与“Node.js + JDK 17 + MySQL 8 手动开发启动”两种交付方式。图~\ref{fig:deploy-arch} 展示了系统容器化部署拓扑结构：Nginx 静态托管 Vue 3 前端并将 \code{/api} 流量同源反代至 Spring Boot 后端容器，后端连接 MySQL 8 容器并挂载 \code{runtime-uploads/} 持久化媒体文件；可选启用 \code{cloudflared} 临时演示隧道或 Cloudflare Pages 云端部署。

\srsfig[0.95\linewidth]{figures/srs/17-system-deployment.png}{系统容器化部署架构、反向代理与本地/云端交付拓扑图}{fig:deploy-arch}

\subsection{部署步骤与运行配置}

\subsubsection{Docker Compose 一键部署流程}
\begin{lstlisting}[language=bash,caption={Docker Compose 一键演示环境初始化与启动命令}]
# 1. 克隆仓库并进入根目录
git clone --branch main https://github.com/Fwlyii/lightweight-takeout-platform.git
cd lightweight-takeout-platform

# 2. 生成本地 .env 并随机签发 64 位 JWT_SECRET（需 zsh 与 OpenSSL）
zsh scripts/ensure-demo-env.sh

# 3. 构建镜像并启动 mysql、migrate、backend、frontend 服务
docker compose -f docker-compose.demo.yml up -d --build
docker compose -f docker-compose.demo.yml ps -a
\end{lstlisting}

首次启动时，\code{mysql} 服务自动创建 \code{elm_v2} 数据库并按序导入 \code{elm_v2.sql} 与 3 份演示种子脚本；随后一次性容器 \code{migrate} 执行 \code{003_order_name_snapshot.sql} 与 \code{005_cart_order_remarks.sql} 后正常退出（\code{Exited (0)}）；待数据库健康检查通过后，\code{backend} 与 \code{frontend} 启动就绪。

\subsubsection{端口映射、访问入口与演示账号}
表~\ref{tab:deploy-endpoints} 列出了默认本地访问入口、服务职责及内建演示账号体系，表~\ref{tab:deploy-env} 列出了外部增强服务的可选环境变量说明。

\begin{longtable}{L{.22\linewidth}L{.30\linewidth}L{.39\linewidth}}
\caption{默认部署端口、访问入口与演示角色说明}\label{tab:deploy-endpoints}\\
\toprule
\textbf{服务/角色} & \textbf{访问地址 / 演示账号} & \textbf{用途与验证说明} \\
\midrule
前端应用（Nginx） & \url{http://localhost:18081/index} & 四端统一入口，已配置同源反向代理至后端容器。 \\
后端服务（Spring Boot） & \url{http://localhost:18080} & 核心 RESTful API 根地址；Swagger 文档见 \code{/swagger-ui/index.html}。 \\
MySQL 8 数据库 & \code{127.0.0.1:13306}（库 \code{elm_v2}） & 持久化存储 23 张业务表及演示数据，支持本地客户端调试。 \\
演示账号矩阵 & \code{demo_user}、\code{demo_merchant}、\code{demo_rider}、\code{demo_rider_candidate}、\code{demo_admin} & 覆盖顾客下单、商家接单、骑手申请/抢单与管理员审核；多角色测试建议开启独立浏览器会话。 \\
\bottomrule
\end{longtable}

\begin{longtable}{L{.28\linewidth}L{.24\linewidth}L{.39\linewidth}}
\caption{关键环境变量与外部服务降级配置表}\label{tab:deploy-env}\\
\toprule
\textbf{环境变量} & \textbf{默认值 / 示例} & \textbf{作用域与缺失降级说明} \\
\midrule
\code{JWT_SECRET} & 由 \code{ensure-demo-env.sh} 生成 & 必填；后端 JWT 令牌签名密钥，禁止提交至代码仓库。 \\
\code{APP_DEMO_ENABLED} & \code{false}（演示环境为 \code{true}） & 控制模拟钱包充值、模拟支付与体验会员接口开关。 \\
\code{DEEPSEEK_API_KEY} & 空（按需注入） & 启用 AI 文字客服外部模型；未配置时自动降级为本地规则与帮助问答。 \\
\code{DASHSCOPE_API_KEY} / \code{BASE_URL} & 空（按需注入） & 启用 \code{qwen-vl-plus} 图片识菜与 \code{qwen3-asr-flash} 语音草稿；未配置返回 503 并禁用入口。 \\
\code{AMAP_WEB_KEY} / \code{AMAP_JS_KEY} / \code{AMAP_SECURITY_CODE} & 空（后端环境变量注入） & 启用高德 POI 搜索代理、逆地理编码与前端选址弹窗；未配置时降级为离线地址填写。 \\
\bottomrule
\end{longtable}

\subsubsection{非 Docker 本地开发启动与常见故障排查}
若需在无 Docker 守护进程的受限主机直接部署调试：可依次启动用户态 MySQL 8.0 实例并导入 \code{elm_v2.sql}、迁移与演示种子脚本，注入 \code{SPRING_DATASOURCE_URL}、\code{JWT_SECRET} 与 \code{APP_DEMO_ENABLED=true} 后使用 OpenJDK 17 运行 \code{elm_bk-0.0.1-SNAPSHOT.jar}；前端通过 \code{VUE_APP_API_BASE_URL} 指向后端地址并运行 \code{npm run serve}（2026 年 11 月 16 日服务器实机截图与全量回归均采用此隔离模式成功验证）。常见部署与运行故障排查指南见表~\ref{tab:deploy-troubleshoot}。

\begin{longtable}{L{.24\linewidth}L{.32\linewidth}L{.35\linewidth}}
\caption{部署与运行阶段常见故障诊断与排查手册}\label{tab:deploy-troubleshoot}\\
\toprule
\textbf{异常现象} & \textbf{常见根因分析} & \textbf{标准排查与处置步骤} \\
\midrule
后端启动报错 \code{JWT_SECRET} 缺失 & 未执行 \code{ensure-demo-env.sh} 或环境变量未透传至容器 & 重新运行环境脚本生成 \code{.env}，检查 \code{docker compose config} 注入结果。 \\
前端接口报 401 或 CORS 拦截 & \code{VUE_APP_API_BASE_URL} 跨端口未走同源反代或密钥变化 & 确认 Nginx \code{/api} 代理指向或统一前端环境变量，清理浏览器旧 \code{localStorage}。 \\
订单页面显示菜名或备注为空 & 既有旧库未执行增量迁移 \code{003} 或 \code{005} & 按序执行 \code{elm_bk/db/migrations/001--006} 幂等 SQL 脚本。 \\
点击高德选址或识菜报不可用 & 未配置 \code{AMAP_*} 或 \code{DASHSCOPE_*} 密钥 & 属预期安全降级行为，可通过文本地址或“智能筛选”正常点餐，配置密钥后重启生效。 \\
\bottomrule
\end{longtable}
